\section{Tata Bahasa dan Derivasi: Definisi Formal}
\label{sec:ch04-grammar}

\begin{definisi}[Tata bahasa]
Tata bahasa \(G\) adalah empat-tupel
\[
G = (V, T, P, S),
\]
dengan \(V\) himpunan berhingga nonterminal, \(T\) himpunan berhingga
terminal, \(P\) himpunan berhingga aturan produksi, dan \(S \in V\)
simbol awal. Himpunan \(V\) dan \(T\) saling lepas.
\end{definisi}

Aturan produksi berbentuk \(\alpha \to \beta\), dengan
\(\alpha \in (V \cup T)^{+}\), \(\beta \in (V \cup T)^{*}\), dan
\(\alpha\) memuat paling sedikit satu nonterminal. Jika
\(\alpha = A \in V\), aturan disebut bebas konteks.

\begin{definisi}[Derivasi]
Jika \(\alpha A \beta\) adalah string dan \(A \to \gamma \in P\),
maka \(\alpha A \beta \Rightarrow \alpha \gamma \beta\) adalah satu
langkah derivasi. Notasi \(\Rightarrow^{*}\) menyatakan nol atau lebih
langkah, dan \(\Rightarrow^{+}\) menyatakan satu atau lebih langkah.
\end{definisi}

Bentuk kalimat adalah string yang dapat diturunkan dari \(S\).
Bahasa yang dibangkitkan \(G\) adalah
\[
L(G) = \{w \in T^{*} \mid S \Rightarrow^{*} w\}.
\]
Dua tata bahasa \(G_1\) dan \(G_2\) ekivalen jika
\(L(G_1) = L(G_2)\).

\subsection{Derivasi Leftmost dan Rightmost}

Derivasi leftmost selalu mengganti nonterminal paling kiri. Derivasi
rightmost mengganti nonterminal paling kanan. Untuk tata bahasa yang
tidak ambigu, keduanya menghasilkan pohon urai yang sama. Ambiguitas
terjadi jika satu string memiliki dua pohon urai berbeda; hal ini
dibahas pada Bab IX.

\begin{contoh}[Derivasi leftmost dan rightmost]
Diberikan tata bahasa
\[
S \to AB, \quad A \to aA \mid a, \quad B \to bB \mid b.
\]
String \(aabb\) dapat diturunkan secara leftmost:
\[
S \Rightarrow AB \Rightarrow aAB \Rightarrow aaB
\Rightarrow aabB \Rightarrow aabb.
\]
Derivasi rightmost:
\[
S \Rightarrow AB \Rightarrow AbB \Rightarrow Abb
\Rightarrow aAbb \Rightarrow aabb.
\]
Kedua derivasi menghasilkan pohon urai yang sama.
\end{contoh}

\subsection{Ekivalensi dan Penyederhanaan}

Tata bahasa dapat disederhanakan dengan menghapus produksi
\(\varepsilon\), produksi unit \(A \to B\), dan simbol yang tidak
berguna. Penyederhanaan ini mempertahankan bahasa dan mempermudah
analisis. Sebagai contoh, produksi unit dapat dihilangkan dengan
menyalin aturan dari nonterminal tujuan.

Penyederhanaan juga membantu mendeteksi masalah desain. Jika sebuah
nonterminal tidak pernah mencapai string terminal, ia tidak berguna.
Jika sebuah nonterminal tidak dapat dicapai dari simbol awal, ia juga
tidak berguna. Kedua kondisi ini sering muncul saat mahasiswa menulis
tata bahasa untuk bahasa yang kompleks.
